\documentclass[11pt,a4paper]{article}
\usepackage{lmodern}
\usepackage[utf8]{inputenc}
\usepackage[T1]{fontenc}
\usepackage{amsmath}
\usepackage{amsfonts}
\usepackage{amssymb}
\usepackage{graphicx}
\usepackage[french]{babel}
\usepackage[notitlepage,fancysections,pagenumber]{sauveterre}

\usepackage{enumitem}

\usepackage{indentfirst}
\usepackage{lipsum}

\usepackage{cases}
\usepackage{multicol}

%formatting titles
\usepackage{titlesec}
\titleformat*{\section}{\centering\Large\bfseries}
\titleformat*{\subsection}{\centering\large\bfseries}

%formatting date
\usepackage[ddmmyyyy]{datetime}
\renewcommand{\dateseparator}{/}

%\usepackage{geometry}
%\geometry{
% a4paper,
% total={170mm,257mm},
% left=20mm,
% top=20mm,
% }

\usepackage{pgf,tikz}
\usepackage{mathrsfs}
\usetikzlibrary{arrows}

\usepackage[export]{adjustbox}

\renewcommand{\notitlepageupperline}{
	\noindent%
	\begin{minipage}{\textwidth}
		\centering
		\begin{minipage}{0.33\textwidth}
			\raisebox{-0.5\height}
			{\color{bleu-sauveterre}%
				\includegraphics[height=0.8\hauteurlogovert]
				{\polylogovert}
				}
		\end{minipage}%
		\begin{minipage}{0.33\textwidth}
			\color{bleu-sauveterre}%
			\centering\polydatesave
			\polyauthorsave
			\hspace{3em}
		\end{minipage}%
		\begin{minipage}{0.33\textwidth}
			\color{bleu-sauveterre}%
			\flushright\raisebox{-0.5\height}
			{\includegraphics[height=0.8\hauteurlogovert]{logo-federal}
			}
		\end{minipage}
	\end{minipage}
}

%\usepackage{fancyhdr}
%\usepackage{lastpage}
%\pagestyle{fancy}
%\fancyhf{}
%%headers and footers
%\lhead{Centre Équestre de Sauveterre}
%\chead{Règlement intérieur}
%\rhead{Version du \today}
%\rfoot{Page \thepage \ sur \pageref{LastPage}}

%\fancypagestyle{firstpage}{%
%  \fancyhf{}
%  \renewcommand\headrulewidth{0pt}
%  \fancyfoot[R]{Page \thepage \ sur \pageref{LastPage}}
%}
%
%\title{\vspace{-50pt}
%\includegraphics[width=0.4\textwidth , left]{logo2-v2.eps}~
%\\[8pt]
%{\huge \textsc{Règlement intérieur}}
%}

\title[Règlement intérieur -- version du \today]{Règlement intérieur}
\subtitle{En date du \today}
\author{}
\date{}

%enumerate environment formatting
\renewcommand{\labelenumi}{\textbf{\arabic{enumi}.}}
\renewcommand{\labelenumii}{\textbf{\arabic{enumi}.\arabic{enumii}.}}
\renewcommand{\labelenumiii}{\textbf{\arabic{enumi}.\arabic{enumii}.\arabic{enumiii}.}}

%change sec and subsec numbering
\renewcommand{\thesection}{Section \Roman{section}.}
\renewcommand{\thesubsection}{Article \arabic{subsection} :}

\renewcommand{\labelenumi}{\arabic{enumi}.}


\renewcommand*\contentsname{Plan}


\usepackage{framed}

%\setlength{\parindent}{0pt}%remove indentation

\begin{document}

\maketitle

%\thispagestyle{firstpage}
%\vspace{-30pt}

Le présent règlement est supposé connu et approuvé par toute personne adhérente au Centre Équestre. Le Centre Équestre s'engage à en adresser une copie à quiconque en ferait la demande.

Le présent règlement abroge et remplace tout règlement précédemment établi.

\subsection*{Préambule}

Ce règlement a comme objectif d'éviter tout accident, incident et malentendu qui pourraient nuire à l'ambiance et à la bonne tenue du Centre Équestre.

Il s'applique dans toute l'enceinte du Centre Équestre (intérieure et extérieure). Il est à respecter par toute personne entrant dans l'enceinte du Centre Équestre : membres, stagiaires, propriétaires, personnels, clientèle.

Le non respect de ce règlement entrainera l'exclusion provisoire ou définitive du Centre Équestre, sans indemnité.

En cas de conflit entre ce qui est énoncé dans le présent règlement et d'autres documents (contrats, affichage,\dots), seul prévaut le contenu du présent règlement. Il en est de même pour ses annexes éventuelles (y compris les \textit{Conditions Générales applicables aux Balades}).

\subsection{Organisation générale}
Toutes les activités du Centre Équestre ainsi que toutes les installations dont il dispose sont placées sous l'autorité des gérants.

\subsection{Fonctionnement}
Pour des raisons de sécurité, tout cavalier doit attendre la présence d'un encadrant du club, ou son accord, avant de s'occuper d'un cheval (hors propriétaire s'occupant de son propre équidé).

Tout cavalier doit impérativement passer au bureau avant le début de la séance pour signaler sa présence à la personne chargée de l'accueil.

Afin d'équilibrer le travail des chevaux et des poneys, les cavaliers sont tenus de signaler, au plus tard la veille, leur éventuelle absence de l'activité.

Il est recommandé aux cavaliers élèves d'être présents une demi-heure avant la reprise.

Les cartes sont strictement nominatives et ne peuvent être revendues ou transférées sans accord de l'exploitant. \par
Les cartes ne sont ni reprises ni remboursées.

Après chaque activité, le matériel doit être remis en place (sauf avis contraire de l'enseignant) :
\begin{itemize}[leftmargin=50pt]
	\item
	Les selles rangées pommeau vers le mur ;
	\item
	Les tapis aux emplacements prévus ;
	\item
	Les étriers remontés ;
	\item
	Les mors lavés ;
	\item
	Les brosses, graisses et accessoires rangés aux places qui leur sont désignées.
\end{itemize}

Les tarifs :
\begin{itemize}[leftmargin=50pt]
\item
Sont affichés au bureau, à l'intérieur et à l'extérieur ;
\item
Sont consultables sur le site internet du Centre Équestre.
\end{itemize}

Les leçons retenues, non décommandées avant 18h00 la veille de la reprise, sont dues.

\subsection{Utilisation des installations}
Toutes les activités de l'établissement ainsi que toutes les installations dont il dispose sont placées sous l'autorité du gérant et du personnel de l'établissement. \par
En période de congés ou d'absence du responsable, les employés restant sur le site sont chargés de veiller au bon fonctionnement du Centre Équestre et d'appliquer les consignes données. \par
Pour la tranquillité et le repos des chevaux, il est demandé de respecter expressément les heures d'ouverture.

Les installations du Centre Équestre -- carrières, manège, écuries, accueil,\dots -- ne sont utilisables que par :
\begin{itemize}[leftmargin=50pt]
\item
Les adhérents à jour de leur cotisation et de leur licence, qui participent à une activité dirigée par des encadrants du Centre Équestre ;
\item
Les cavaliers \og de passage \fg{} qui participent à une activité encadrée par le Centre Équestre ;
\item
Les propriétaires qui utilisent seuls leur cheval (\textit{cf.} article 10).
\end{itemize}

Toute autre cas d'utilisation des installations est formellement interdit. Une éventuelle dérogation requiert obligatoirement et préalablement l'accord écrit du Centre Équestre.

Seuls sont autorisés à pratiquer le saut d'obstacles en autonomie les propriétaires titulaires du Galop 7. Les autres en ont l'interdiction formelle.

Chaque membre doit veiller au respect des installations et à la propreté générale des lieux.

L'accès aux écuries, pour le public, est interdit en dehors des horaires d'ouverture du Centre Équestre. Pour les propriétaires, hors cas particuliers (cheval malade, autorisation préalable de la direction,\dots), l'accès aux écuries et aux installations est \textbf{interdit avant 7h30 et au-delà de 21h30}.

Des restrictions ponctuelles ou permanentes concernant l'utilisation des installations et aires de travail peuvent être instaurées par la direction et seront notifiées aux cavaliers et propriétaires.

Les véhicules doivent être garés sur le parking ou des emplacements en bordure de route prévus à cet effet.

\subsection{Interdiction de fumer}
Il est strictement interdit de fumer dans l'enceinte du Centre Équestre, à l'exception des endroits prévus à cet effet.

Les mégots devront impérativement jetés dans les cendriers prévus à cet effet.

\subsection{Observations et suggestions par les membres}
Un cahier est tenu à la disposition des adhérents, à l'accueil, afin que ceux-ci puissent y consigner les observations et suggestions qu'ils désirent formuler concernant le fonctionnement du Centre Équestre.

\subsection{Discipline}
Durant les activités, et en particulier à l'intérieur des locaux ou installations, les adhérents doivent observer une obéissance complète aux encadrants, en particulier concernant les consignes de sécurité.

En tout lieu et en toute circonstance, les adhérents sont tenus d'observer une attitude respectueuse empreinte de correction et de savoir-vivre vis-à-vis du personnel encadrant et des autres salariés.

En retour et dans les mêmes circonstances, le personnel encadrant et les salariés sont tenus d'observer une attitude respectueuse empreinte de correction et de savoir-vivre vis-à-vis des adhérents.

Aucune manifestation discourtoise envers le Centre Équestre, ses structures, ses administrateurs, son personnel ou ses adhérents n'est admise par oral. (\textit{cf.} article 7).

Toute personne fréquentant le Centre Équestre se doit d'observer les règles de sécurité internes (ne pas courir, faire de gestes brusques,...) dans les écuries ou les lieux de travail.

Les jeux (ballons, boomerangs,...) sont interdits dans l'enceinte du Club.\par
Les chiens extérieurs à ceux appartenant au Centre Équestre doivent être tenus en laisse.\par
Les jeux des enfants dans les bottes de paille sont également interdits.

D'une manière générale, les enfants ne doivent pas jouer sur le site sans être surveillés de leurs parents. Ils sont sous leur responsabilité.

Pendant les reprises, seul l'enseignant et ses aides sont habilités à se trouver sur le terrain d'évolution.

Il est interdit de sortir ou de déplacer les obstacles des carrières ou manèges sans l'autorisation du responsable.

\subsection{Réclamations}
Tout adhérent désireux de présenter une réclamation qu'il estime motivée et justifiée concernant le Centre Équestre peut opérer de l'une des manières suivantes :
\begin{itemize}[leftmargin=50pt]
\item
S'adresser directement à l'un des encadrants de façon confidentielle ;
\item
S'adresser directement aux gérants de l'exploitation par oral ou leur adresser un courrier ;
\item
Consigner sa réclamation sur le cahier prévu à l'article 5.
\end{itemize}
Tout réclamation présentée sous l'une des formes ainsi définies doit recevoir une réponse sous 8 jours.

\subsection{Sanctions}
Toute attitude répréhensible d'un adhérent et en particulier toute inobservation des statuts ou du règlement intérieur expose celui qui en est responsable à des sanctions qui peuvent être de trois ordres~:
\begin{enumerate}
\item
La mise à pied prononcée par le Comité Exécutif pour une durée ne pouvant excéder un mois. L'adhérent qui est mis à pied ne peut, pendant la durée de la sanction, ni monter un cheval appartenant au Centre Équestre, ni utiliser les terrains d'évolution, manège ou carrières.
\item
L'exclusion temporaire ou suspension, prononcée par le Conseil d'Administration pour une durée ne pouvant excéder une année. L'adhérent qui est exclus temporairement n'a plus accès aux locaux et installations du Centre Équestre, et ne peut, pendant la durée de la sanction, participer à aucune des activités -- publiques ou privées -- ni assister aux assemblées générales.
\item
L'exclusion définitive.
\end{enumerate}
Tout adhérent faisant l'objet d'une sanction ne peut prétendre à aucun remboursement des sommes déjà payées par lui se rapportant aux activités dont la sanction le prive.

\textit{\textbf{Le respect du cheval étant une règle absolue, tout mauvais traitement entraîne l'exclusion immédiate.}}

\subsection{Tenue}
Pour monter à cheval, tant à l'intérieur qu'à l'extérieur, les adhérents du Centre Équestre doivent adopter une tenue vestimentaire correcte et conforme aux usages traditionnels de l'équitation française.

Le port du casque est obligatoire. Il doit être porté afin de constituer une protection effective pour le cavalier et être conforme à la norme imposée par la règlementation en vigueur.

\subsection{Assurances}
Les adhérents sont obligatoirement assurés pour leur responsabilité civile dans le cadre de l'assurance du Centre Équestre, durant le temps d'activité équestre. Il leur appartient de prendre connaissance au bureau de l'étendue et des limites de garantie qui leur sont ainsi accordées.

Aucun adhérent ne peut participer aux activités du Centre Équestre ou utiliser ses installations s'il n'a pas acquitté sa cotisation pour l'année en cours et s'il n'est pas titulaire d'une licence de la FFE de l'année en cours.

La responsabilité du Centre Équestre est dégagée dans le cas d'un accident provoqué par une inobservation du règlement intérieur.

\subsection{Prise en charge des enfants mineurs}
Les cavaliers mineurs ne sont sous la responsabilité du Centre Équestre que durant la durée de leur reprise (ou de leur promenade) et durant le temps de préparation de l'équidé et le retour à l'écurie (ou en prairie), soit \textit{\textbf{une demi-heure avant}} et \textit{\textbf{une demi-heure après}}.

En dehors des heures des activités ainsi définies, même sur le site du Centre Équestre, \textit{\textbf{les mineurs sont sous la responsabilité de leurs parents ou de leur tuteur légal}}.

\subsection{Balades et cartes cadeau}

Tout cavalier, y compris non adhérent se doit de prendre connaissance des \textit{Conditions Générales applicables aux Balades}, décrivant leurs modalités et leur fonctionnement, ainsi que les conditions de remboursement.
\\

Les \og cartes cadeau \fg{} sont des coupons prépayés émis par le Centre Équestre pouvant être déclinés pour tout type de prestation, et édités après paiement de la prestation voulue ou à titre promotionnel. Elles ont vocation à être offertes, et sont d'une validité variant de 6 mois à un an en fonction de la prestation. Une carte cadeau ne peut être ni échangée, ni remboursée, ni voir sa validité prolongée (sauf cas de force majeure, sur accord de la direction). Elle devra être présentée le jour de l'activité,  avant l'accès à la prestation.

\subsection{Règlement des prestations et remboursement}

Les prestations doivent être réglées à l'accueil du Centre Équestre, en arrivant, le jour de la prestation. Toute balade doit être réglée avant son déroulement, et le Centre Équestre se réserve le droit d'accepter ou non au départ d'une balade quelqu'un ne l'ayant pas réglée selon ces modalités.

Toutes les prestations proposées par le Centre Équestre peuvent être réglées \textit{via} les moyens suivant :
\begin{itemize}[leftmargin=50pt]
  \item
  Espèces de toute valeur (\textit{NB} : la direction se réserve le droit de refuser les coupures de 200 et 500 euros) ;
  \item
  chèque bancaire (à l'ordre de \texttt{SARL Centre Équestre de Sauveterre}) ;
  \item
  virement bancaire (demander un RIB à l'accueil si nécessaire) ;
  \item
  chèques vacances ;
  \item
  coupons sport ;
	\item carte cadeau émise par le Centre Équestre, en cours de validité.
\end{itemize}

La direction se réserve le droit de refuser tout moyen de paiement ne figurant pas dans cette liste.

Pour les cartes de cours et forfaits, un paiement en plusieurs fois est possible en accord avec les responsables.

\paragraph{Modalités de remboursement :\\}
Concernant les modalités de remboursement des balades, se référer aux \textit{Conditions Générales applicables aux Balades}.

Pour toutes les autres prestations, le remboursement est décidé au cas par cas entre la direction et les intéressés.


\subsection{Propriétaires de chevaux}
Les propriétaires auront pris connaissance des \textit{Conditions Générales de prise en pension}, annexe II du présent règlement.


\subsection{Armoires de rangement}

Afin de ranger et protéger les affaires, des armoires de rangement peuvent être louées. Deux tailles sont proposées (petite ou grande) et elles sont toutes fermées à clé ou cadenas. Ces \og casiers \fg{} peuvent être loués par tout cavalier du club, propriétaire ou non, seul ou à plusieurs, dans la limite des disponibilités. De plus, tout locataire de casier devra indiquer son nom sur son casier, en utilisant l'étiquette prévue à cet effet.

Le paiement est mensuel, et \textit{\textbf{tout casier non réglé pendant une période de 6 mois sera libéré et les affaires rangées à l'accueil jusqu'à régularisation du paiement}}. Évidemment, tout casier libéré peut être attribué à un autre cavalier en faisant la demande.

Sont mis à disposition des cavaliers, dans chaque casier, au minimum 1 porte-selle et 1 porte-filet pour les petits casiers, et 2 de chaque dans les grands. Les cavaliers sont libres de ranger ce qu'ils veulent dans leurs, casiers, cependant \textit{\textbf{merci de ne pas y ranger de nourriture, ou alors la garder dans des boîtes hermétiques}} pour des raisons évidentes d'hygiène.

\subsection{Autorisation de reproduction d'images et cession de droits}
Conformément aux dispositions relatives au droit à l'image et au droit d'auteur, les adhérents autorisent le Centre Équestre ainsi que son Association à capter, utiliser ou diffuser des images photo ou vidéo prises dans le cadre des activités du Club, sur les supports suivants :
\begin{itemize}
\item
Site internet du Club ;
\item
Journal d'information du Club ;
\item
Documents de présentation et promotion des activités de la SARL Centre Équestre de Sauveterre;
\item
Journaux de la presse quotidienne locale ;
\item
Journaux spécialisés des sports équestres.
\end{itemize}
Dans les conditions d'utilisation définies ci-dessus, les adhérents renoncent à tout recours ou toute indemnité contre le Centre Équestre ou son Association.

\subsection{Application}
En signant leur adhésion à l'établissement, les membres reconnaissent formellement avoir pris connaissance du règlement intérieur et en accepter toutes les dispositions.

\vfill

Fait aux Sables d'Olonne, le \today

\vspace{2cm}

\begin{flushright}
\rule{6cm}{0.5pt}

{Stéphane MARAIS, Directrice \kern 9pt}

\end{flushright}

\end{document}
